\chapter{Algorithms and Pseudocode}
\label{app:algorithms}

\renewcommand{\thealgocf}{\thechapter.\arabic{algocf}}
\setcounter{algocf}{0}

This appendix collects compact pseudocode for the attacks and training procedures used in the report. The algorithms use symbolic budgets, step counts, and learning rates; concrete configuration values are defined in Chapter~\ref{chap:experimental_setup}.

\begin{algorithm}[H]
\caption{FGSM-RS}
\DontPrintSemicolon
\KwIn{Model \(f_\theta\), input \(x\), label \(y\), loss \(\ell\), budget \(\varepsilon\), step size \(\alpha\)}
\KwOut{Adversarial input \(x^{\mathrm{adv}}\)}
Sample \(\delta_0 \in \mathcal{S}_\infty(\varepsilon)\)\;
Set \(x_0 \leftarrow \Pi_{[0,1]^d}(x+\delta_0)\)\;
Compute \(g \leftarrow \nabla_x \ell(f_\theta(x_0),y)\)\;
Set \(x^{\mathrm{adv}} \leftarrow \Pi_{(x+\mathcal{S}_\infty(\varepsilon))\cap[0,1]^d}(x_0+\alpha\,\mathrm{sign}(g))\)\;
\Return{\(x^{\mathrm{adv}}\)}
\end{algorithm}

\begin{algorithm}[H]
\caption{PGD Attack}
\DontPrintSemicolon
\KwIn{Model \(f_\theta\), input \(x\), label \(y\), loss \(\ell\), set \(\mathcal{S}_p(\varepsilon)\), step size \(\alpha\), steps \(T\)}
\KwOut{Adversarial input \(x_T^{\mathrm{adv}}\)}
Initialize \(x_0^{\mathrm{adv}}\in (x+\mathcal{S}_p(\varepsilon))\cap[0,1]^d\)\;
\For{\(t=0,\ldots,T-1\)}{
    Compute \(g_t \leftarrow \nabla_x \ell(f_\theta(x_t^{\mathrm{adv}}),y)\)\;
    Take a normalized ascent step using the \(\ell_p\) geometry\;
    Project \(x_{t+1}^{\mathrm{adv}}\) onto \((x+\mathcal{S}_p(\varepsilon))\cap[0,1]^d\)\;
}
\Return{\(x_T^{\mathrm{adv}}\)}
\end{algorithm}

\begin{algorithm}[H]
\caption{CW-\(\ell_2\) Attack}
\DontPrintSemicolon
\KwIn{Model \(f_\theta\), input \(x\), label \(y\), margin loss \(g\), trade-off \(c\), optimizer steps \(T\)}
\KwOut{Adversarial input \(x^{\mathrm{adv}}\)}
Initialize perturbation variable \(\delta\)\;
\For{\(t=1,\ldots,T\)}{
    Minimize \(\|\delta\|_2^2+c\,g(x+\delta)\) by gradient-based optimization\;
    Clip or transform \(x+\delta\) so the candidate remains in \([0,1]^d\)\;
    Keep the best successful adversarial candidate found so far\;
}
\Return{best successful candidate, or the final candidate if none succeeds}
\end{algorithm}

\begin{algorithm}[H]
\caption{DeepFool-\(\ell_2\) Attack}
\DontPrintSemicolon
\KwIn{Model \(f_\theta\), input \(x\), current class \(y\), steps \(T\)}
\KwOut{Adversarial input \(x^{\mathrm{adv}}\)}
Set \(x_0^{\mathrm{adv}}\leftarrow x\)\;
\For{\(t=0,\ldots,T-1\)}{
    \If{\(\arg\max_c f_\theta(x_t^{\mathrm{adv}})_c \neq y\)}{
        \Return{\(x_t^{\mathrm{adv}}\)}
    }
    Linearize the class scores around \(x_t^{\mathrm{adv}}\)\;
    Estimate the closest class boundary under the local linear model\;
    Move \(x_t^{\mathrm{adv}}\) by the corresponding minimal \(\ell_2\) perturbation\;
}
\Return{\(x_T^{\mathrm{adv}}\)}
\end{algorithm}

\begin{algorithm}[H]
\caption{DDN-\(\ell_2\) Attack}
\DontPrintSemicolon
\KwIn{Model \(f_\theta\), input \(x\), label \(y\), loss \(\ell\), initial radius \(\rho_0\), step size \(\alpha\), steps \(T\)}
\KwOut{Adversarial input \(x^{\mathrm{adv}}\)}
Initialize \(\delta_0\) and radius \(\rho_0\)\;
\For{\(t=0,\ldots,T-1\)}{
    Compute \(g_t \leftarrow \nabla_x \ell(f_\theta(x+\delta_t),y)\)\;
    Update perturbation direction using the normalized gradient\;
    \If{\(x+\delta_t\) is adversarial}{
        Decrease the radius for the next candidate\;
    }
    \Else{
        Increase the radius for the next candidate\;
    }
    Normalize the direction to the updated radius and clip to \([0,1]^d\)\;
}
\Return{best successful candidate found during the search}
\end{algorithm}

\begin{algorithm}[H]
\caption{MI-FGSM}
\DontPrintSemicolon
\KwIn{Model \(f_\theta\), input \(x\), label \(y\), loss \(\ell\), budget \(\varepsilon\), step size \(\alpha\), momentum \(\mu\), steps \(T\)}
\KwOut{Adversarial input \(x_T^{\mathrm{adv}}\)}
Initialize \(x_0^{\mathrm{adv}}\leftarrow x\) and \(g_0\leftarrow 0\)\;
\For{\(t=0,\ldots,T-1\)}{
    Compute \(r_t \leftarrow \nabla_x \ell(f_\theta(x_t^{\mathrm{adv}}),y)\)\;
    Set \(g_{t+1}\leftarrow \mu g_t + r_t/\|r_t\|_1\)\;
    Set \(x_{t+1}^{\mathrm{adv}}\leftarrow \Pi_{(x+\mathcal{S}_\infty(\varepsilon))\cap[0,1]^d}(x_t^{\mathrm{adv}}+\alpha\,\mathrm{sign}(g_{t+1}))\)\;
}
\Return{\(x_T^{\mathrm{adv}}\)}
\end{algorithm}

\begin{algorithm}[H]
\caption{TRADES-style Projected Gradient Descent (TPGD)}
\DontPrintSemicolon
\KwIn{Model \(f_\theta\), input \(x\), budget \(\varepsilon\), step size \(\alpha\), steps \(T\)}
\KwOut{Adversarial input \(x_T^{\mathrm{adv}}\)}
Initialize \(x_0^{\mathrm{adv}}\in (x+\mathcal{S}_\infty(\varepsilon))\cap[0,1]^d\)\;
Compute clean prediction \(p_\theta(\cdot\mid x)\)\;
\For{\(t=0,\ldots,T-1\)}{
    Compute \(L_t \leftarrow \operatorname{KL}(p_\theta(\cdot\mid x)\,\|\,p_\theta(\cdot\mid x_t^{\mathrm{adv}}))\)\;
    Ascend \(\nabla_x L_t\) and project onto \((x+\mathcal{S}_\infty(\varepsilon))\cap[0,1]^d\)\;
}
\Return{\(x_T^{\mathrm{adv}}\)}
\end{algorithm}

\begin{algorithm}[H]
\caption{AutoAttack Evaluation Wrapper}
\DontPrintSemicolon
\KwIn{Model \(f_\theta\), evaluation set \(\mathcal{D}_{\mathrm{test}}\), AutoAttack component suite \(\mathcal{A}_{\mathrm{AA}}\)}
\KwOut{AutoAttack robust accuracy}
Initialize success indicator \(s_i\leftarrow 1\) for each test example\;
\For{each component attack \(A\in\mathcal{A}_{\mathrm{AA}}\)}{
    Generate adversarial candidates \(A(x_i,y_i,f_\theta)\)\;
    Set \(s_i\leftarrow 0\) for examples misclassified by the component attack\;
}
Compute robust accuracy \(\frac{1}{|\mathcal{D}_{\mathrm{test}}|}\sum_i s_i\)\;
\Return{robust accuracy}
\end{algorithm}

\begin{algorithm}[H]
\caption{Multi-AT Training}
\DontPrintSemicolon
\KwIn{Training set \(\mathcal{D}_{\mathrm{train}}\), model \(f_\theta\), source attack set \(\mathcal{E}_{\mathrm{src}}\), loss \(\ell\), learning rate \(\eta_\theta\)}
\KwOut{Trained model \(f_\theta^\ast\)}
\For{each training epoch}{
    \For{each minibatch \(\mathcal{B}\)}{
        Assign examples in \(\mathcal{B}\) uniformly across attacks in \(\mathcal{E}_{\mathrm{src}}\)\;
        Generate adversarial examples using the assigned source attack\;
        Compute the uniform mean adversarial loss over the minibatch\;
        Update \(\theta\) using the model optimizer and learning rate \(\eta_\theta\)\;
    }
}
\Return{\(f_\theta^\ast\)}
\end{algorithm}

\begin{algorithm}[H]
\caption{AttackDRO Training}
\DontPrintSemicolon
\KwIn{Training set \(\mathcal{D}_{\mathrm{train}}\), model \(f_\theta\), source attack set \(\mathcal{E}_{\mathrm{src}}\), attack weights \(q\), learning rates \(\eta_\theta,\eta_q\)}
\KwOut{Trained model \(f_\theta^\ast\)}
Initialize \(q\) uniformly over source attack identities\;
\For{each training epoch}{
    \For{each minibatch}{
        Generate adversarial examples using the source attacks\;
        Estimate loss or difficulty for each source attack group\;
        Compute the attack-weighted Group DRO loss using \(q\)\;
        Update \(\theta\) with the weighted loss\;
        Update \(q\) using exponentiated-gradient reweighting over attack groups\;
    }
}
\Return{\(f_\theta^\ast\)}
\end{algorithm}

\begin{algorithm}[H]
\caption{Cluster-DRO and AttackDRO++ Training Overview}
\DontPrintSemicolon
\KwIn{Training set \(\mathcal{D}_{\mathrm{train}}\), model \(f_\theta\), source attacks \(\mathcal{E}_{\mathrm{src}}\), cluster count \(K\), anchor weight \(\lambda_{\mathrm{DRO}}\), group weights \(q\)}
\KwOut{Trained model \(f_\theta^\ast\)}
Initialize feature bank, clusterer, and uniform group weights \(q\)\;
\For{each training epoch}{
    Refresh latent clusters from the feature bank when the refresh condition is met\;
    \For{each minibatch}{
        Generate source adversarial examples as in Multi-AT\;
        Extract representations and, for AttackDRO++, augmented clustering features with gradient fingerprints\;
        Assign each adversarial example to a latent cluster\;
        Compute the uniform adversarial loss over the minibatch\;
        Compute the Cluster-DRO loss using cluster weights \(q\)\;
        Combine the uniform anchor and Cluster-DRO terms using \(\lambda_{\mathrm{DRO}}\)\;
        Update model parameters and then update \(q\) from cluster-level difficulty estimates\;
        Add detached augmented features to the feature bank\;
    }
}
\Return{\(f_\theta^\ast\)}
\end{algorithm}
